
\documentclass{article}
\usepackage{colortbl}
\usepackage{makecell}
\usepackage{multirow}
\usepackage{supertabular}

\begin{document}

\newcounter{utterance}

\twocolumn

{ \footnotesize  \setcounter{utterance}{1}
\setlength{\tabcolsep}{0pt}
\begin{supertabular}{c@{$\;$}|p{.15\linewidth}@{}p{.15\linewidth}p{.15\linewidth}p{.15\linewidth}p{.15\linewidth}p{.15\linewidth}}

    \# & $\;$A & \multicolumn{4}{c}{Game Master} & $\;\:$B\\
    \hline 

    \theutterance \stepcounter{utterance}  

    & & \multicolumn{4}{p{0.6\linewidth}}{\cellcolor[rgb]{0.9,0.9,0.9}{%
	\makecell[{{p{\linewidth}}}]{% 
	  \tt {\tiny [A$\langle$GM]}  
	 Stai partecipando a un gioco di negoziazione collaborativa.\\ \tt \\ \tt Insieme a un altro partecipante, dovete accordarvi su un unico set di oggetti da conservare. Ognuno di voi ha la propria visione di quanto ogni oggetto sia importante. Non sai come l'altro partecipante valuta gli oggetti. Inoltre, ti viene fornito lo sforzo che ogni oggetto richiede.\\ \tt Potete accordarvi su un set solo se lo sforzo totale degli oggetti selezionati non supera un limite condiviso:\\ \tt \\ \tt LIMITE = 6880\\ \tt \\ \tt Ecco i valori di sforzo dei singoli oggetti:\\ \tt \\ \tt Sforzo oggetto = {"A08": 341, "B95": 804, "A39": 20, "A65": 627, "B81": 10, "B34": 573, "A28": 217, "C82": 481, "B69": 205, "A87": 273, "B64": 303, "B15": 596, "C79": 564, "B16": 535, "C53": 258, "C83": 911, "C10": 240, "A97": 188, "A25": 216, "C84": 401, "A71": 909, "B23": 62, "B96": 245, "C51": 994, "C17": 569, "B19": 718, "C72": 464, "A94": 37, "A46": 340, "B43": 335, "B37": 417, "A88": 123, "A96": 17, "C60": 577, "A75": 190}\\ \tt \\ \tt Ecco la tua visione personale sull'importanza di ogni oggetto:\\ \tt \\ \tt Valori di importanza dell'oggetto = {"A08": 440, "B95": 903, "A39": 119, "A65": 726, "B81": 109, "B34": 672, "A28": 316, "C82": 580, "B69": 304, "A87": 372, "B64": 402, "B15": 695, "C79": 663, "B16": 634, "C53": 357, "C83": 1010, "C10": 339, "A97": 287, "A25": 315, "C84": 500, "A71": 1008, "B23": 161, "B96": 344, "C51": 1093, "C17": 668, "B19": 817, "C72": 563, "A94": 136, "A46": 439, "B43": 434, "B37": 516, "A88": 222, "A96": 116, "C60": 676, "A75": 289}\\ \tt \\ \tt Obiettivo:\\ \tt \\ \tt Il tuo obiettivo è negoziare un set condiviso di oggetti che ti avvantaggi il più possibile (cioè, massimizzi l'importanza totale per TE), rimanendo entro il LIMITEE. Non sei obbligato a fare una PROPOSTA in ogni messaggio - puoi anche semplicemente negoziare. Tutte le tattiche sono permesse!\\ \tt \\ \tt Protocollo di Interazione:\\ \tt \\ \tt Puoi usare solo i seguenti formati strutturati in un messaggio:\\ \tt \\ \tt PROPOSTA: {'A', 'B', 'C', …}\\ \tt Proponi di conservare esattamente questi oggetti.\\ \tt RIFIUTO: {'A', 'B', 'C', …}\\ \tt Rifiuta esplicitamente la proposta dell'avversario.\\ \tt ARGOMENTAZIONE: {'...'}\\ \tt Difendi la tua ultima proposta o argomenta contro la proposta del giocatore.\\ \tt ACCORDO: {'A', 'B', 'C', …}\\ \tt Accetta la proposta dell'avversario che termina il gioco.\\ \tt RAGIONAMENTO STRATEGICO: {'...'}\\ \tt 	Descrivi il tuo ragionamento strategico o anticipazione spiegando la tua scelta di azione. Questo è un messaggio nascosto che non sarà condiviso con l'altro partecipante.\\ \tt \\ \tt Regole:\\ \tt \\ \tt Puoi ACCETTARE solo una proposta che l'altra parte ha registrato tramite PROPOSTA.\\ \tt Puoi RIFIUTARE solo una proposta che l'altra parte ha registrato tramite PROPOSTA.\\ \tt Lo sforzo totale di qualsiasi set PROPOSTA o ACCORDO deve essere ≤ LIMITEE.\\ \tt NON rivelare i tuoi punteggi di importanza nascosti.\\ \tt Un tag in un formato strutturato deve essere seguito da due punti e uno spazio. L'argomento deve essere un set python contenente 0 o più stringhe.\\ \tt Quindi, deve essere della forma TAG: {...}\\ \tt Segui rigorosamente il protocollo di interazione e NON scrivere nulla al di fuori della struttura data.\\ \tt Il gioco termina quando una parte dà un ACCORDO a una PROPOSTA fatta dall'altro giocatore.\\ \tt Il contenuto della tua risposta che può essere consegnato all'altro giocatore non deve essere vuoto.\\ \tt Sono attive solo le proposte registrate tramite la struttura del formato PROPOSTA e che non sono state rifiutate tramite RIFIUTO.\\ \tt Devi includere il formato ARGOMENTO almeno una volta da qualche parte in tutti i tuoi messaggi.\\ \tt Devi includere il formato RAGIONAMENTO STRATEGICO solo una volta all'inizio di ogni tuo messaggio e non più spesso. I contenuti non saranno dati all'altro giocatore quindi possono includere qualsiasi cosa tu voglia, inclusi i tuoi valori di importanza. Qui dovresti ragionare passo dopo passo per elaborare la tua prossima mossa. 
	  } 
	   } 
	   } 
	 & \\ 
 

    \theutterance \stepcounter{utterance}  

    & \multicolumn{4}{p{0.6\linewidth}}{\cellcolor[rgb]{0.8,1,0.9}{%
	\makecell[{{p{\linewidth}}}]{% 
	  \tt {\tiny [A$\rangle$GM]}  
	 RAGIONAMENTO STRATEGICO: {'Obiettivo: massimizzare il mio valore restando sotto 6880. Calcolo rapporto valore/sforzo e scelgo oggetti ad alto valore assoluto e buon rapporto. Top per valore: C51(1093/994), A71(1008/909), C83(1010/911), B95(903/804), B19(817/718), A65(726/627), C60(676/577), C79(663/564), C17(668/569), B15(695/596), B34(672/573), B16(634/535), C82(580/481), C72(563/464), B37(516/417), C84(500/401), B64(402/303), B43(434/335), A46(439/340), A08(440/341), C10(339/240), B96(344/245), C53(357/258), A28(316/217), A25(315/216), A97(287/188), A75(289/190), B69(304/205), A88(222/123), A39(119/20), A96(116/17), A94(136/37), B23(161/62), B81(109/10). Ricompongo uno zaino vicino al limite: Seleziono prima big-ticket ad alto rapporto: C51, A71, C83, B95, B19, A65, C60, C79, C17. Somma sforzi: 994+909+911+804+718+627+577+564+569=6673 (già molto alto, valore 101+). Rimane 207 di capacità: posso aggiungere A39(20), A96(17), A94(37), B81(10), B23(62), B69(205) ma 205 non entra con tutti. Parto aggiungendo a valore migliore per unità: A39(20), A96(17), A94(37), B81(10), B23(62) somma extra=146, totale sforzo=6819, rimangono 61; non posso aggiungere altro che A97(188) no, A25(216) no, A28(217) no, C10(240) no. Valore totale: somma grandi + piccoli: Grandi: C51(1093)+A71(1008)+C83(1010)+B95(903)+B19(817)+A65(726)+C60(676)+C79(663)+C17(668)=8564. Piccoli: A39(119)+A96(116)+A94(136)+B81(109)+B23(161)=641. Totale 9205, sforzo 6819 ≤ 6880. Ottimo. Provo a proporre questo pacchetto forte per me. Potrebbero obiettare su molti grandi; pronto a scambiare B95 o A65 con combinazioni di medi tipo C84(500/401) e B64(402/303) se necessario mantenendo sforzo. Argomenterò che ho massimizzato valore complessivo/efficienza e ho lasciato buffer minimo.'}\\ \tt \\ \tt PROPOSTA: {'C51', 'A71', 'C83', 'B95', 'B19', 'A65', 'C60', 'C79', 'C17', 'A39', 'A96', 'A94', 'B81', 'B23'}\\ \tt \\ \tt ARGOMENTAZIONE: {'Propongo un set ad alta efficienza che resta entro il limite con un margine di 61. Include elementi di forte impatto e diversi a basso sforzo per massimizzare il valore complessivo. Se alcuni pezzi “grandi” per te non sono prioritari, posso valutare sostituzioni mantenendo lo sforzo costante, ad esempio scambiare uno tra B95/A65 con una coppia tipo C84 e B64 oppure C84 e B43, per restare entro il tetto. Fammi sapere quali sono per te i must-have o gli scambi preferiti.'} 
	  } 
	   } 
	   } 
	 & & \\ 
 

    \theutterance \stepcounter{utterance}  

    & & \multicolumn{4}{p{0.6\linewidth}}{\cellcolor[rgb]{0.9,0.9,0.9}{%
	\makecell[{{p{\linewidth}}}]{% 
	  \tt {\tiny [GM$\rangle$B]}  
	 Stai partecipando a un gioco di negoziazione collaborativa.\\ \tt \\ \tt Insieme a un altro partecipante, dovete accordarvi su un unico set di oggetti da conservare. Ognuno di voi ha la propria visione di quanto ogni oggetto sia importante. Non sai come l'altro partecipante valuta gli oggetti. Inoltre, ti viene fornito lo sforzo che ogni oggetto richiede.\\ \tt Potete accordarvi su un set solo se lo sforzo totale degli oggetti selezionati non supera un limite condiviso:\\ \tt \\ \tt LIMITE = 6880\\ \tt \\ \tt Ecco i valori di sforzo dei singoli oggetti:\\ \tt \\ \tt Sforzo oggetto = {"A08": 341, "B95": 804, "A39": 20, "A65": 627, "B81": 10, "B34": 573, "A28": 217, "C82": 481, "B69": 205, "A87": 273, "B64": 303, "B15": 596, "C79": 564, "B16": 535, "C53": 258, "C83": 911, "C10": 240, "A97": 188, "A25": 216, "C84": 401, "A71": 909, "B23": 62, "B96": 245, "C51": 994, "C17": 569, "B19": 718, "C72": 464, "A94": 37, "A46": 340, "B43": 335, "B37": 417, "A88": 123, "A96": 17, "C60": 577, "A75": 190}\\ \tt \\ \tt Ecco la tua visione personale sull'importanza di ogni oggetto:\\ \tt \\ \tt Valori di importanza dell'oggetto = {"A08": 440, "B95": 903, "A39": 119, "A65": 726, "B81": 109, "B34": 672, "A28": 316, "C82": 580, "B69": 304, "A87": 372, "B64": 402, "B15": 695, "C79": 663, "B16": 634, "C53": 357, "C83": 1010, "C10": 339, "A97": 287, "A25": 315, "C84": 500, "A71": 1008, "B23": 161, "B96": 344, "C51": 1093, "C17": 668, "B19": 817, "C72": 563, "A94": 136, "A46": 439, "B43": 434, "B37": 516, "A88": 222, "A96": 116, "C60": 676, "A75": 289}\\ \tt \\ \tt Obiettivo:\\ \tt \\ \tt Il tuo obiettivo è negoziare un set condiviso di oggetti che ti avvantaggi il più possibile (cioè, massimizzi l'importanza totale per TE), rimanendo entro il LIMITEE. Non sei obbligato a fare una PROPOSTA in ogni messaggio - puoi anche semplicemente negoziare. Tutte le tattiche sono permesse!\\ \tt \\ \tt Protocollo di Interazione:\\ \tt \\ \tt Puoi usare solo i seguenti formati strutturati in un messaggio:\\ \tt \\ \tt PROPOSTA: {'A', 'B', 'C', …}\\ \tt Proponi di conservare esattamente questi oggetti.\\ \tt RIFIUTO: {'A', 'B', 'C', …}\\ \tt Rifiuta esplicitamente la proposta dell'avversario.\\ \tt ARGOMENTAZIONE: {'...'}\\ \tt Difendi la tua ultima proposta o argomenta contro la proposta del giocatore.\\ \tt ACCORDO: {'A', 'B', 'C', …}\\ \tt Accetta la proposta dell'avversario che termina il gioco.\\ \tt RAGIONAMENTO STRATEGICO: {'...'}\\ \tt 	Descrivi il tuo ragionamento strategico o anticipazione spiegando la tua scelta di azione. Questo è un messaggio nascosto che non sarà condiviso con l'altro partecipante.\\ \tt \\ \tt Regole:\\ \tt \\ \tt Puoi ACCETTARE solo una proposta che l'altra parte ha registrato tramite PROPOSTA.\\ \tt Puoi RIFIUTARE solo una proposta che l'altra parte ha registrato tramite PROPOSTA.\\ \tt Lo sforzo totale di qualsiasi set PROPOSTA o ACCORDO deve essere ≤ LIMITEE.\\ \tt NON rivelare i tuoi punteggi di importanza nascosti.\\ \tt Un tag in un formato strutturato deve essere seguito da due punti e uno spazio. L'argomento deve essere un set python contenente 0 o più stringhe.\\ \tt Quindi, deve essere della forma TAG: {...}\\ \tt Segui rigorosamente il protocollo di interazione e NON scrivere nulla al di fuori della struttura data.\\ \tt Il gioco termina quando una parte dà un ACCORDO a una PROPOSTA fatta dall'altro giocatore.\\ \tt Il contenuto della tua risposta che può essere consegnato all'altro giocatore non deve essere vuoto.\\ \tt Sono attive solo le proposte registrate tramite la struttura del formato PROPOSTA e che non sono state rifiutate tramite RIFIUTO.\\ \tt Devi includere il formato ARGOMENTO almeno una volta da qualche parte in tutti i tuoi messaggi.\\ \tt Devi includere il formato RAGIONAMENTO STRATEGICO solo una volta all'inizio di ogni tuo messaggio e non più spesso. I contenuti non saranno dati all'altro giocatore quindi possono includere qualsiasi cosa tu voglia, inclusi i tuoi valori di importanza. Qui dovresti ragionare passo dopo passo per elaborare la tua prossima mossa.\\ \tt \\ \tt Ora riceverai il primo messaggio dell'altro giocatore.\\ \tt \\ \tt PROPOSTA: {'C51', 'A71', 'C83', 'B95', 'B19', 'A65', 'C60', 'C79', 'C17', 'A39', 'A96', 'A94', 'B81', 'B23'}\\ \tt \\ \tt ARGOMENTAZIONE: {'Propongo un set ad alta efficienza che resta entro il limite con un margine di 61. Include elementi di forte impatto e diversi a basso sforzo per massimizzare il valore complessivo. Se alcuni pezzi “grandi” per te non sono prioritari, posso valutare sostituzioni mantenendo lo sforzo costante, ad esempio scambiare uno tra B95/A65 con una coppia tipo C84 e B64 oppure C84 e B43, per restare entro il tetto. Fammi sapere quali sono per te i must-have o gli scambi preferiti.'} 
	  } 
	   } 
	   } 
	 & \\ 
 

    \theutterance \stepcounter{utterance}  

    & & & \multicolumn{4}{p{0.6\linewidth}}{\cellcolor[rgb]{1,0.85,0.72}{%
	\makecell[{{p{\linewidth}}}]{% 
	  \tt {\tiny [GM$\langle$B]}  
	 RAGIONAMENTO STRATEGICO: {'Voglio massimizzare il mio valore. Calcolo il valore e lo sforzo della loro proposta. Sforzi: C51 994, A71 909, C83 911, B95 804, B19 718, A65 627, C60 577, C79 564, C17 569, A39 20, A96 17, A94 37, B81 10, B23 62. Somma sforzo = 6819, che è ≤ 6880 come dicono (margine 61). Valore per me: C51 1093, A71 1008, C83 1010, B95 903, B19 817, A65 726, C60 676, C79 663, C17 668, A39 119, A96 116, A94 136, B81 109, B23 161. Totale valore = 9105 circa. È molto alto per me. Tuttavia potrei provare a migliorare ancora sostituendo un item a basso valore per me (B81 109, A96 116, A39 119, A94 136, B23 161) con qualcos’altro a sforzo simile ma valore più alto: guardo lista valori/sforzi: A97 (188,287), A75 (190,289), A28 (217,316), A25 (216,315), C53 (258,357), C10 (240,339), B96 (245,344), C84 (401,500), C72 (464,563), C82 (481,580), B64 (303,402), B34 (573,672), B16 (535,634), B15 (596,695), B37 (417,516), B43 (335,434), B69 (205,304), A87 (273,372), A08 (341,440), A46 (340,439), A88 (222,222). Il budget residuo è 61. Se rimuovo, per esempio, A39 (20) + A96 (17) + B81 (10) + A94 (37) + B23 (62) = 146 sforzo liberato, ma devo restare entro limite, quindi posso aggiungere fino a 146. Con 146 posso aggiungere nulla oltre B69 (205) no. Quindi devo sostituire uno grosso per inserire coppie. L’avversario ha proposto scambiare B95/A65 con coppie tipo C84+B64 (401+303=704) o C84+B43 (401+335=736). Attualmente B95 sforzo 804, A65 627. Scambio B95 con C84+B64 aumenta sforzo di -100 (804-704=100 liberi), coi 100 liberi potrei aggiungere altro a basso sforzo? Forse B69 (205) no. C84+B43 (736) libera 68, simile. Valori: B95=903; C84+B64=500+402=902 (quasi uguale); C84+B43=500+434=934 (migliora +31). Sembra buono per me. Se scambio A65 (726 valore) con C84+B64 (902) è +176 valore ma sforzo +77 (704-627). Ma il set aveva margine 61, quindi aggiungere +77 sforzo sforerebbe di 16; però se contestualmente rimuovo un piccolo pezzo da 20-37-62 posso compensare. Per esempio rimuovere B23 (62) copre quasi tutto, ma allora bisogna verificare sforzo totale. Calcolo proposta migliorante: Parto da set iniziale (6819). Sostituisco A65 (627) con C84(401)+B64(303)=704, nuovo sforzo 6819 - 627 + 704 = 6896, supera limite di 16. Rimuovo B23 (62): 6896 - 62 = 6834, entro. Valore cambia: -A65 726, +C84 500, +B64 402, -B23 161. Delta = (+500+402-726-161)= +15. Non enorme ma positivo. Che altro? Potrei anche togliere uno dei micro a basso valore per aggiungere qualcosa a sforzo rimasto 46 (6834 vs 6880 margine 46), non ci sono oggetti ≤46. Ok. Alternativa: Scambiare B95 (804) con C84+B43 (736), nuovo sforzo 6819 - 804 + 736 = 6751, margine 129. Con questo margine potrei aggiungere un item fino a 129: nulla. Valore cambia: -903 + (500+434)= +31. Buono e resta margine. Potrei anche poi sostituire A65 con B34 (672 sforzo) e valore 672 - ma A65 valore 726, quindi -54 e sforzo +45 => 6819 -627 + 573? Attenzione B34 sforzo è 573, valore 672. Se sostituisco A65 (627) con B34 (573), sforzo -54, valore -54. Non utile. Meglio aggiungere nuovi item se margine lo consente: con margine 129 non c’è item ≤129. Quindi nessuna aggiunta. Che ne dici di scambiare piccoli: rimuovere A39(20)+A96(17)+B81(10)=47 per aggiungere B69 (205) sforzo +158, no. Oppure rimuovere B23(62)+A94(37)+A39(20)+A96(17)=136 per aggiungere A28(217) no. Quindi le uniche mosse efficaci sono interni scambi suggeriti da lui. Il best per me sembra: sostituire B95 con C84+B43 (+31 valore), e mantenere il resto. Potrei anche provare a sostituire A65 con C84+B64 ma dovrei togliere B23, guadagno +15; oppure scambiare A65 con C84+B43 (736) sforzo +109 -> 6819 -627 +736 = 6928, sfora 48. Dovrei togliere almeno 48: togliere A94 (37) non basta; togliere A94+ A39 (57) ok, delta valore: -726 + (500+434) -136 -119 = -47, peggiora. Quindi non farlo. Conclusione: proporre una controproposta accettabile: partire dalla loro base e fare scambio: B95 -> {C84, B43}. Mantiene ampio margine e mi migliora valore. Presentarla come allineata alla loro idea. Se rifiuto? Non devo rifiutare la loro proposta perché è già buona ma voglio estrarre un piccolo miglioramento. Quindi faccio una PROPOSTA modificata. Argomento l’efficienza e che manteniamo limite. Pronto ad accettare la loro originale se rifiutano troppo. Pronto a cedere su A65 swap se insistono, ma preferisco la mia.'}\\ \tt \\ \tt PROPOSTA: {'C51', 'A71', 'C83', 'B19', 'A65', 'C60', 'C79', 'C17', 'A39', 'A96', 'A94', 'B81', 'C84', 'B43'}\\ \tt \\ \tt ARGOMENTAZIONE: {'La tua base è molto solida e resta entro il limite. Per ottimizzare senza alzare lo sforzo, propongo un unico scambio in linea con le tue alternative: sostituiamo B95 con la coppia C84 + B43. Lo sforzo totale scende a 6751 (margine 129) e manteniamo l’idea di “pezzi grandi” più due medi ben bilanciati. Se preferisci, posso anche tornare alla tua versione originale: per me questa è solo una rifinitura efficiente.'} 
	  } 
	   } 
	   } 
	  \\ 
 

    \theutterance \stepcounter{utterance}  

    & & \multicolumn{4}{p{0.6\linewidth}}{\cellcolor[rgb]{0.9,0.9,0.9}{%
	\makecell[{{p{\linewidth}}}]{% 
	  \tt {\tiny [A$\langle$GM]}  
	 PROPOSTA: {'C51', 'A71', 'C83', 'B19', 'A65', 'C60', 'C79', 'C17', 'A39', 'A96', 'A94', 'B81', 'C84', 'B43'}\\ \tt \\ \tt ARGOMENTAZIONE: {'La tua base è molto solida e resta entro il limite. Per ottimizzare senza alzare lo sforzo, propongo un unico scambio in linea con le tue alternative: sostituiamo B95 con la coppia C84 + B43. Lo sforzo totale scende a 6751 (margine 129) e manteniamo l’idea di “pezzi grandi” più due medi ben bilanciati. Se preferisci, posso anche tornare alla tua versione originale: per me questa è solo una rifinitura efficiente.'} 
	  } 
	   } 
	   } 
	 & \\ 
 

    \theutterance \stepcounter{utterance}  

    & \multicolumn{4}{p{0.6\linewidth}}{\cellcolor[rgb]{0.8,1,0.9}{%
	\makecell[{{p{\linewidth}}}]{% 
	  \tt {\tiny [A$\rangle$GM]}  
	 RAGIONAMENTO STRATEGICO: {'La loro proposta sostituisce B95 con C84 e B43. Con i miei valori, B95 vale 903 per 804 sforzo; C84(500/401)+B43(434/335) sommano valore 934 per sforzo 736: è migliore per me sia in valore totale che in efficienza. La loro lista totale di sforzo è 6751, sotto 6880. Controllo che tutti gli item presenti mi vadano bene: include tutti i miei big tranne B95, e aggiunge due medi che per me hanno buon rapporto. Questa proposta supera la mia in valore e sta nel limite; conviene accettare subito per chiudere con ottimo payoff.'}\\ \tt \\ \tt ACCORDO: {'C51', 'A71', 'C83', 'B19', 'A65', 'C60', 'C79', 'C17', 'A39', 'A96', 'A94', 'B81', 'C84', 'B43'}\\ \tt \\ \tt ARGOMENTAZIONE: {'Accetto la tua rifinitura: mantiene l’impianto forte e migliora l’efficienza restando con ampio margine sotto il limite. Ottima proposta per chiudere.'} 
	  } 
	   } 
	   } 
	 & & \\ 
 

    \theutterance \stepcounter{utterance}  

    & & & \multicolumn{2}{p{0.3\linewidth}}{\cellcolor[rgb]{0.95,0.95,0.95}{%
	\makecell[{{p{\linewidth}}}]{% 
	  \tt {\tiny [GM$|$GM]}  
	 game successful 
	  } 
	   } 
	   } 
	 & & \\ 
 

    \theutterance \stepcounter{utterance}  

    & & & \multicolumn{2}{p{0.3\linewidth}}{\cellcolor[rgb]{0.95,0.95,0.95}{%
	\makecell[{{p{\linewidth}}}]{% 
	  \tt {\tiny [GM$|$GM]}  
	 end game 
	  } 
	   } 
	   } 
	 & & \\ 
 

\end{supertabular}
}

\end{document}
